\zlabel{triptych:concise:commentary:start}
\section{The Propers: Detailed Commentary}

\subsection{What does belonging give, and what does it ask?}

Augustine's search for the garment begins with goods that Christians
already possess. Baptism, attendance at the altar, fasting, knowledge
and exceptional spiritual gifts are serious goods. Yet good and bad
recipients can share them. Their possession therefore cannot by itself
explain why one guest honours the feast while another stands condemned.
Augustine asks what makes these gifts profitable to the person who
receives them. Charity answers that question because it concerns the
life animating participation, not simply an additional object acquired
alongside the others.\footnote{Augustine, Sermon 90.5--6; NPNF I.6,
Sermon XL in the English sequence, Benedictine XC.}

Chrysostom's account of the new person begins at a different point.
He explains how the same human being can become new. Paul's clothing
language does not replace one human nature with another. A former way
of life is changed; the person is not imprisoned in an unalterable
identity as a wrongdoer. At the same time, the new person's creation
in righteousness says more than a change of name. God's action makes
the commanded life possible. The baptismal garment is to remain visible
in conduct.\footnote{Chrysostom, \work{Homilies on Ephesians} XIII,
on 4:22--24. The appointment begins at 4:23; verse 22 supplies context.}

These explanations concern two questions about a received gift.
Augustine asks why its reception can be fruitless; Chrysostom asks
how its recipient can change. The first refuses presumption, the second
despair. Neither makes divine gift a reward for prior perfection, and
neither treats a religious identity as an excuse to remain unchanged.
Their moral demands arise within belonging to Christ. The thief who
must now work and share is the person called to receive, not someone
who must buy an invitation with the proceeds of his improvement.

The Introit's psalm brings the same question to the act of hearing.
Its larger history remembers gifts accepted without steadfast trust.
Augustine addresses Christians through that history and also recognizes
faithful recipients of grace among ancient Israel. The warning cannot
be turned into a contrast between a faithless people then and an
automatically faithful Church now. It is a question put to the present
hearer: what does the gift make of the one who receives it?

Augustine and Bellarmine disagree more narrowly over the opening voice.
Augustine hears God speaking through the prophet; Bellarmine identifies
David as the immediate speaker, reasoning from the subsequent account
of what the fathers have told. They agree that hearing must submit to
divine instruction, but their identification of the speaker is not the
same. Neither gives a biblical address to the composed \textit{Salus
populi} antiphon. The psalm verse's summons follows the antiphon's
promise; the two texts retain their distinct identities.\footnote{Augustine,
\work{Enarrationes} 77.1--3; Bellarmine, \work{Psalms} 77:1--8.}

The Collect makes the appropriate response a request for freedom.
Its body and mind are to serve God together. Schuster describes their
ordered service: bodily strength has a good end, and the soul is raised
to God through grace. This is neither flight from the body nor the
promise that illness will never impede activity. Its freedom is for
service. Augustine's fruitful guest and Chrysostom's renewed person
both need what the congregation asks for here: a life whose powers
can answer the giver.\footnote{Schuster, \work{Sacramentary} III,
pp.~171--172.}

\subsection{Why is charity the garment of this wedding?}

Gregory reaches charity through the identity of the Bridegroom.
Christ came in love to unite the Church to himself. The garment is
therefore appropriate to this wedding because it answers the love
celebrated there. The uncharitable guest refuses the very good that
gives the feast its meaning. Augustine reasons from the insufficiency
of possessions shared by good and bad; Gregory reasons from Christ's
loving union. Their conclusions agree, but the two arguments do
different work. One searches the recipient; the other discloses the
gift to which he must respond.\footnote{Gregory, \work{Homilies on the
Gospels} II.38.3, 9; Augustine, Sermon 90.5--6.}

The allegorical meaning thus belongs to Christ and his body, not to
a social expectation about costly clothes. Matthew does not establish
a universal custom of hosts providing wedding garments. Gregory's
Christology also sets the terms of the nuptial image. The Incarnation
does not join two persons in marriage: Christ is one Person in two
natures. The union of Christ and the Church unfolds the love of that
one Lord. An image drawn from a human wedding cannot substitute its
own arrangements for the mystery it helps to explain.

Gregory's weaving image extends this union into the twofold command
of love. A woven garment rests upon upper and lower supports; charity
requires love of God and of neighbour. Neglect of a neighbour's need
does not become harmless through professions of devotion, and service
of a neighbour cannot dispense with love of God. Augustine subjects
even almsgiving to a corresponding question of intention. A deed can
be impressive while its actor seeks admiration. The neighbour's
actual good and the giver's love must meet.\footnote{Gregory,
\work{Gospel Homilies} 38.10; Augustine, Sermon 90.6. Gregory's
image concerns the upper and lower supports of weaving, not two threads.}

Bellarmine's account of the Gradual sharpens this distinction in
prayer. Incense rising before God images a right direction of the
will. A wish to be admired draws prayer away from that direction.
Faith, hope, love and humility belong to acceptable prayer, which
rests upon Christ's priesthood. The visible gesture of lifted hands
therefore invites examination of the invisible aim; it does not
make external worship dispensable. Augustine places the evening
sacrifice in Christ's voluntary Passion and the praying Church in
her Head. Her access to God is received through his self-gift.
Her love is a participation in that gift, not an offering that repairs
some insufficiency in it.\footnote{Bellarmine, \work{Psalms} 140:1--2;
Augustine, \work{Enarrationes} 140.2--4. The verse gives a sacrificial
image, not a rubric fixing this Mass's hour or prescribing a gesture
at the Gradual.}

The charity reading consequently gives restored action its direction.
Hands can work energetically and lips can speak eloquently while the
heart remains turned toward itself. The renewal reading supplies
the converse test: a love claiming to unite the person to Christ
must reach what those hands and lips actually do. The literal acts
and the allegorical union illuminate one another without becoming
interchangeable. Neither eloquence nor bodily effort is itself the
garment; neither can be sealed off from the love the garment signifies.

\subsection{How does love reach the difficult neighbour?}

Chrysostom explains truthfulness through the body's mutual dependence.
An eye that hid danger from the foot would harm the body to which
it belongs. The neighbour is someone who must act upon what he is
told. A lie therefore does more than stain its speaker: it damages
the trust by which another member can live safely. Paul does not
merely recommend politeness. He identifies the common life that
makes falsehood destructive to the liar as well as to the deceived.
Truthful speech serves another's good; it is not permission for
indiscriminate disclosure.\footnote{Chrysostom, \work{Ephesians}
XIV, on 4:25.}

Augustine and Gregory extend the demand beyond friends who readily
repay kindness. Augustine asks the Christian to desire the evil
person's healing rather than destruction. Christ and Stephen praying
for persecutors give this desire its pattern. Gregory likewise
tests charity by adverse hatred: the friend is loved in God, the
enemy for God's sake. The demand is not to declare the evil good.
It is to seek the person's freedom from the evil, refusing to make
his destruction one's own good.\footnote{Augustine, Sermon 90.7--10;
Gregory, \work{Gospel Homilies} 38.11.}

Chrysostom's treatment of anger explains a practical obstacle to
that desire. Injury is rehearsed, resentment grows, and hostility
becomes the judge of ambiguous words and deeds. A slander then
finds a willing listener because it fits an already cherished
picture of the enemy. Sunset interrupts this growth. Chrysostom
describes how the night's quieter hours can feed what the day
has begun. Delay is not neutral; it can furnish anger with a more
secure dwelling.\footnote{Chrysostom, \work{Ephesians} XIV,
on 4:26--27.}

The immediate task is therefore more precise than manufacturing
calm by a deadline. A wrong may still require correction, and its
consequences may not be settled by evening. Yet the injured person
can stop nursing revenge and refuse the inward pleasure of making
an enemy irredeemable. Augustine's prayer for healing supplies
the end toward which Chrysostom's interruption of resentment can
move. One explains the breadth of charity; the other explains a
process by which the heart loses that breadth.

The Offertory supplies confidence for such exposure. The singer
remains in tribulation and acknowledges hostile anger. Augustine
distinguishes the life God secures from the temporal safety an
enemy can threaten. Charity need not wait for a world in which
no enemy can inflict injury. Nor does it make injury unreal.
The saving hand belongs to God; hatred does not have to become
the worshipper's means of ultimate self-preservation.\footnote{Augustine,
\work{Enarrationes} 137.10--13.}

The Gradual and Alleluia give this healed speech further work.
The psalm surrounding the Gradual asks that the mouth be guarded.
Bellarmine explains the need for both understanding and strength:
to know when to speak and to speak or remain silent accordingly.
Even cooperation is grace's gift. The Alleluia then opens the
mouth in praise and proclamation. Bellarmine sees the lover's
wish that others know the beloved; Augustine hears the evangelists
in the announcement to the nations. Restraint from harmful words
is fulfilled in a tongue able to invite and give thanks.\footnote{Bellarmine,
\work{Psalms} 140:3 and 104:1--2; Augustine, \work{Enarrationes}
104.1--3. Psalm 140:3 is context for the appointed verse 2.}

\subsection{When do working hands become generous hands?}

Paul ends the appointed Epistle with a person in need. The former
thief must work at what is good so that he can give. Chrysostom
insists on both changes: taking must cease, and positive help must
follow. Labour considered merely as effort is no sufficient good;
the thief also exerts himself. Its object and purpose matter.
Nor is the command exhausted when the worker can support his own
ambitions. Someone else's need enters the end for which the hands
labour.\footnote{Chrysostom, \work{Ephesians} XIV, on 4:28.}

This particular demand exposes what a general exhortation to charity
could leave untouched. A person may praise generosity while always
arranging his work and possessions around himself. The Epistle does
not let love remain an inward sentiment detached from the disposition
of what one has. Conversely, Augustine's account of intention prevents
the visible gift from becoming its own complete vindication. Giving
for glory can still curve inward toward the actor. The renewed hands
must reach the neighbour's actual need in a love that seeks his good.

Gregory reads the bound hands and feet of the judged guest as an
exposure of culpable inaction. Hands withheld alms and feet neglected
the sick; their refusal in life anticipated their binding in judgment.
That is his moral exposition of the scene, not the literal cause
Matthew supplies for the bonds. It concerns refusing help that one
could give. Disability, poverty and inability to work are not the
absent garment. Paul puts the needy person at the receiving end of
the command, not outside the feast for failing to produce enough.\footnote{Gregory, \work{Gospel Homilies} 38.13; compare Augustine,
Sermon 90.6.}

The Gospel also shows why productive occupation cannot be a person's
final good. Those who prefer field or merchandise can remain busy
while refusing the feast. The invitation calls them into a relation
with the Son, and work finds its place within that relation. The
worker is first a guest. Efficiency cannot replace the communion
for which his powers are renewed.

Prayer and action therefore need one another without becoming the
same act. The Gradual's lifted hands ask acceptance; the Epistle's
hands supply need; the Offertory's divine hand saves. The Eucharistic
gift is not simply a name for human labour's product. Christ gives
the grace through which human action can become charity. Schuster's
exposition of the Secret locates the congregation precisely here:
the saving gift is real, and the participants ask for its saving
fruit in themselves. What they offer does not exempt them from
receiving or from conversion.\footnote{Schuster, \work{Sacramentary}
III, p.~173; Augustine, \work{Enarrationes} 140.2--4.}

\subsection{Can obedience be both commanded and received?}

The Communion's two verses answer in their very succession.
First comes God's command, then the worshipper's wish for ways
directed toward it. Augustine explains \textit{nimis} as great
diligence rather than excessive obedience, and \textit{utinam}
as petition. The command is not weakened by the request for help.
The request acknowledges that knowing the good is not identical
with possessing the strength to do it.\footnote{Augustine,
\work{Enarrationes} 118, opening Aleph exposition, 4--5.}

Bellarmine gives this dependence a distinction between beginning
and growth. Initial justification is not the payment for law
observed without grace. Yet the justified person really grows
through obedience enabled by grace. Divine help does not make
action unreal, and the reality of action does not turn it into
independent self-justification. Even the psalm's later firm
resolution is joined to a request not to be forsaken.\footnote{Bellarmine,
\work{Psalms} 118:4--8; verse 8 supplies context, while the
Communion appoints verses 4--5.}

Schuster hears delight in the same wish. Augustine's dependence
and Schuster's joy are complementary, not competing explanations
that require one author to be discarded. Dependence answers the
temptation to boast; delight answers the thought that God's service
must remain alien to what a person loves. Bellarmine's growth
answers the temptation to make dependence a reason for doing
nothing. The speaker can desire the right path, need guidance,
and genuinely advance upon it.

Aquinas's exposition of renewal gives the inner change several
expressions. The ``spirit of your mind'' can refer to the Holy
Spirit dwelling in the mind, the rational spirit, or the mind
made spiritual. His account converges on holiness in the heart,
truth in the mouth and justice in action, with Christ as the
beginning of newness. Chrysostom's particular commands show those
powers in use. Neither commentator's account makes conversion
merely external, and Aquinas's alternatives prevent one phrase
from being forced into an artificially single explanation.\footnote{Aquinas, \work{Super Ephesios} 4, lectio 7, Dessain
1857, vol.~II, p.~329; paraphrase of the inspected Latin.
Schuster, \work{Sacramentary} III, pp.~173--174.}

The Postcommunion asks that God's medicine reach this inward
life. Release from perversity leads to adherence to the commands.
Schuster describes grace healing desire, not simply supplying
information about the good. A person may know what should be
done while continuing to wish for its opposite. The prayer asks
for a freedom in which desire and action become ordered to God.
The Collect's freedom for service returns as the healing of
what has made service difficult from within.

Here the two readings meet most closely. Charity needs growth
in the recipient; ordinary action needs a healed source in the
will. The one who prays for an enemy's healing must also accept
the diagnosis contained in the congregation's own prayer.
The Postcommunion is plural. Those who have heard the guest
examined now ask that they themselves be freed and kept faithful.
The medicine belongs to those who receive it, not merely to the
people they judge to need it.

\subsection{What hope survives the king's judgment?}

Gregory distinguishes charity imperfectly possessed from charity
wholly absent. The distinction is decisive for the anxious guest.
Weak love needs nourishment; it is not thereby no love at all.
Augustine likewise speaks of charity being born, nourished and
increased. Neither author permits indifference, but neither
requires the believer to regard every struggle as evidence of
exclusion. The garment has a life of growth before it reaches
perfection.\footnote{Gregory, \work{Gospel Homilies} 38.12;
Augustine, Sermon 90.2--3, 6.}

The mixed hall is therefore a present reality, not a portrait
of evil eternally sharing the final good. Augustine distinguishes
the feast now from its perfected fulfilment. Gregory favours a
relation between Matthew's and Luke's banquets in which the
present mixed assembly differs from the final banquet from
which no admitted guest departs, while allowing another account
of their relation. Their theological destination is compatible;
Gregory does not require one exhaustive harmonization of the
narratives.\footnote{Augustine, Sermon 90.1--4; Gregory,
\work{Gospel Homilies} 38.1--3, 7.}

The judgment remains real. The single expelled guest represents
a class in Augustine's exposition, not a fraction from which
to calculate a salvation rate. Gregory turns the unknown outcome
of one's own perseverance toward humility: good beginnings can
be abandoned and a misdirected life can repent. Neither the
violence of the parable nor the king's judgment makes human
retaliation a Christian duty. The hearer stands under examination
and is called to conversion.\footnote{Augustine, Sermon 90.4;
Gregory, \work{Gospel Homilies} 38.13--14.}

The anagogical accent differs within the two whole-formulary
readings. The garment reading looks toward a feast in which
charity is perfected and evil excluded. The renewal reading
looks toward human life made whole: its labour and trial are
gathered into a destiny beyond their present visible success.
Augustine's Passion interpretation of the evening sacrifice and
Schuster's resurrection morning give that hope its Christological
centre. These are two emphases within one end, not competing
destinies.\footnote{Augustine, \work{Enarrationes} 140.3;
Schuster, \work{Sacramentary} III, pp.~172--174.}

The Postcommunion's request for lasting adherence is consequently
both practical and eternal. It concerns the next truthful word,
the refusal to feed resentment and the gift another person needs;
it also asks that the life begun through grace endure. The
believer does not need to claim perfection to hope for the feast.
He needs to receive its healing and let the love of the Son
become a life.
